\documentclass[12pt,letterpaper]{article}

% ---------------------------------------------------------------
% Configuración general
% ---------------------------------------------------------------
\usepackage{fontspec}
\usepackage{polyglossia}
\setdefaultlanguage{spanish}
\setmainfont{Liberation Serif}
\setmonofont{Liberation Mono}

\usepackage{geometry}
\geometry{margin=2.5cm}
\usepackage{xcolor}
\usepackage{graphicx}
\usepackage{tikz}
\usetikzlibrary{arrows.meta,positioning,fit,calc}
\usepackage{listings}
\usepackage{booktabs}
\usepackage{tabularx}
\usepackage{array}
\usepackage{enumitem}
\usepackage{ragged2e}
\usepackage{float}
\usepackage{hyperref}
\usepackage{cite}

\definecolor{codegray}{gray}{0.95}
\definecolor{codegreen}{rgb}{0.0,0.45,0.0}
\definecolor{codeblue}{rgb}{0.1,0.1,0.6}
\definecolor{diagramblue}{RGB}{218,232,252}
\definecolor{diagramgreen}{RGB}{226,239,218}
\definecolor{diagramorange}{RGB}{252,228,214}

\lstset{
    backgroundcolor=\color{codegray},
    basicstyle=\ttfamily\small,
    keywordstyle=\color{codeblue}\bfseries,
    commentstyle=\color{codegreen}\itshape,
    stringstyle=\color{red!70!black},
    numbers=left,
    numberstyle=\tiny\color{gray},
    numbersep=8pt,
    breaklines=true,
    showstringspaces=false,
    frame=single,
    rulecolor=\color{gray!50},
    tabsize=2,
    extendedchars=true
}

\hypersetup{
    colorlinks=true,
    linkcolor=blue!50!black,
    citecolor=blue!50!black,
    urlcolor=blue!50!black,
    pdftitle={Práctica 4: Gestión concurrente de pacientes y vacunas},
    pdfauthor={Equipo de Programación Concurrente y Paralela}
}

\newcolumntype{Y}{>{\RaggedRight\arraybackslash}X}

\title{Práctica 4. Gestión concurrente de pacientes y vacunas}
\author{Islas Gómez Luis Antonio \and Ramírez Candia Alfredo \and Pérez Flores Julio César \and Navarro Soto Mario Alberto}
\date{\today}

\begin{document}

% ---------------------------------------------------------------
% Portada basada en el reporte de referencia del equipo
% ---------------------------------------------------------------
\begin{titlepage}
    \centering
    \vspace*{0.5cm}

    {\Large\bfseries BENEMÉRITA UNIVERSIDAD AUTÓNOMA DE PUEBLA\par}
    \vspace{0.4cm}
    {\large Facultad de Ciencias de la Computación\par}
    \vspace{1.2cm}

    {\large Programación Concurrente y Paralela\par}

    \vfill

    {\Huge\bfseries Práctica 4\par}
    \vspace{0.4cm}
    {\LARGE Gestión concurrente de pacientes y vacunas\par}
    \vspace{0.3cm}
    {\large (Segundo parcial)\par}

    \vfill

    {\large\bfseries Equipo:\par}
    \vspace{0.3cm}
    Islas Gómez Luis Antonio --- 202372860\par
    Ramírez Candia Alfredo --- 202375947\par
    Pérez Flores Julio César --- 202247445\par
    Navarro Soto Mario Alberto --- 202264602\par

    \vfill

    {\large Docente: Hilda Castillo Zacatelco\par}
    \vspace{0.6cm}
    {\large \today\par}
\end{titlepage}

\tableofcontents
\newpage

% ---------------------------------------------------------------
\section{Resumen}

Este reporte documenta el análisis, diseño e implementación de una aplicación
para administrar pacientes pediátricos y sus vacunas. La información persistente
se almacena en un archivo de texto, mientras que las operaciones de registro y
actualización se ejecutan sobre una estructura en memoria. La solución emplea
hilos de C++ y exclusión mutua para permitir operaciones concurrentes sin
corromper el vector compartido de pacientes.

Se describen las decisiones tomadas para resolver cada requisito, el modelo de
dependencias entre tareas, la representación conceptual mediante bloques
\textsc{cobegin}/\textsc{coend} y las condiciones que justifican la
sincronización. El código fuente se incluye únicamente en los apéndices; la
explicación principal utiliza modelos, tablas, diagramas y pseudocódigo.

\noindent\textbf{Palabras clave:} concurrencia, hilos, exclusión mutua, vacunas,
persistencia, grafo de dependencias, COBEGIN--COEND.

% ---------------------------------------------------------------
\section{Introducción y objetivo}

La práctica plantea el desarrollo de un sistema para una clínica de medicina
general que atiende pacientes pediátricos. Cada registro contiene el nombre del
paciente, su fecha de nacimiento, las vacunas recibidas y las vacunas próximas a
aplicarse. Al abrir la clínica, los registros deben cargarse desde un archivo;
al cerrarla, los cambios deben persistirse nuevamente.

El objetivo de esta implementación no es únicamente resolver las consultas de
la clínica, sino estudiar qué partes del flujo pueden ejecutarse de forma
concurrente y cuáles deben protegerse. La estrategia adoptada combina un
vector compartido en memoria, tareas independientes representadas por hilos y
un mutex que serializa cada operación que modifica o inspecciona dicho vector.
Esta decisión sigue el principio de que el paralelismo debe aplicarse sobre
tareas independientes, mientras que el acceso a un estado común debe tener una
sección crítica bien definida \cite{foster1995,andrewsmultithreaded,herlihy2020}.

\subsection{Objetivos específicos}

\begin{itemize}[leftmargin=1.2cm]
    \item Modelar pacientes, vacunas y fechas mediante clases separadas.
    \item Validar fechas y descartar registros de entrada que no cumplan el
    formato establecido.
    \item Implementar las operaciones solicitadas sobre pacientes y vacunas.
    \item Ejecutar altas y actualizaciones independientes mediante hilos.
    \item Proteger el estado compartido con exclusión mutua.
    \item Representar el flujo paralelo mediante un grafo de dependencias y
    pseudocódigo \textsc{cobegin}/\textsc{coend}.
    \item Verificar la persistencia y la compilación reproducible del sistema.
\end{itemize}

% ---------------------------------------------------------------
\section{Descripción del problema}

El archivo de entrada utiliza una línea por paciente y cuatro campos separados
por punto y coma:

\begin{enumerate}[leftmargin=1.2cm]
    \item nombre del paciente;
    \item fecha de nacimiento;
    \item lista de vacunas recibidas, representadas como pares nombre--fecha;
    \item lista de vacunas próximas, también representadas como pares
    nombre--fecha.
\end{enumerate}

Los campos vacíos son válidos para expresar que un paciente todavía no tiene
vacunas recibidas o próximas. En cambio, una fecha con formato incorrecto, una
fecha imposible o un nombre vacío no puede incorporarse al modelo interno. La
separación entre análisis sintáctico y validación semántica evita que un dato
mal formado provoque una excepción no controlada durante la carga.

La especificación exige las siguientes operaciones: registrar pacientes,
registrar vacunas aplicadas, agregar vacunas próximas, eliminar pacientes,
modificar datos personales, modificar o eliminar vacunas, realizar búsquedas,
ordenar pacientes, consultar rangos de edad y producir estadísticas temporales.
Además, el archivo debe mantenerse como punto de partida de las ejecuciones
posteriores y conservar el orden lógico de los registros.

% ---------------------------------------------------------------
\section{Estado del arte}

\subsection{Concurrencia, paralelismo y tareas}

La concurrencia describe la composición de varias actividades que pueden estar
en progreso durante un mismo intervalo; el paralelismo ocurre cuando esas
actividades se ejecutan físicamente al mismo tiempo en distintos núcleos. La
diferencia es importante en esta práctica: crear varios hilos expresa
concurrencia, pero la ganancia de tiempo solo aparece si el hardware y la carga
permiten ejecutar las tareas en paralelo \cite{andrewsmultithreaded,herlihy2020}.

Foster propone analizar un problema paralelo mediante partición, comunicación,
aglomeración y mapeo \cite{foster1995}. Aplicado a la clínica, la partición
consiste en separar altas y actualizaciones por paciente; la comunicación es el
acceso al objeto `Clinica`; la aglomeración evita crear una tarea para cada
operación trivial, y el mapeo asigna cada grupo de operaciones a un hilo.

\subsection{Memoria compartida y exclusión mutua}

La solución utiliza memoria compartida dentro del proceso: todos los hilos
pueden acceder al objeto que contiene el vector de pacientes. Esta estrategia
es directa, pero introduce condiciones de carrera si dos hilos leen y escriben
el vector al mismo tiempo. El mutex representa la exclusión mutua: solo un hilo
puede entrar a la sección crítica protegida en un instante dado.

El uso de `std::lock\_guard` implementa adquisición y liberación con alcance
automático. Esta propiedad es relevante porque evita dejar el mutex bloqueado
cuando una función retorna anticipadamente o se produce una excepción. El
estándar de C++ define estas primitivas como parte de su biblioteca de soporte
para concurrencia \cite{iso14882,cppreferenceconcurrency}.

\subsection{COBEGIN--COEND y grafos de dependencias}

La notación \textsc{cobegin}/\textsc{coend} expresa un bloque de actividades que
comienzan de manera concurrente y cuyo final constituye una barrera. Aunque el
programa se implementa con `std::thread` y `join`, la correspondencia conceptual
es directa: la creación de hilos equivale al comienzo de las ramas y cada
operación `join` participa en el `COEND`.

Los grafos de dependencias permiten distinguir precedencias obligatorias de
actividades independientes. Una arista significa que la tarea destino no puede
comenzar correctamente hasta que la tarea origen produzca su resultado. En la
práctica, la carga del archivo precede a todas las operaciones; el alta produce
los identificadores que requieren las actualizaciones; y el `COEND` precede a
las consultas finales. Este análisis evita paralelizar operaciones que todavía
no tienen los datos necesarios \cite{dijkstra1968,hoare1978}.

% ---------------------------------------------------------------
\section{Análisis y diseño de la solución}

\subsection{Modelo de datos}

El diseño separa las responsabilidades en cuatro entidades principales:

\begin{itemize}[leftmargin=1.2cm]
    \item \textbf{Paciente:} mantiene un identificador, nombre, fecha de
    nacimiento y dos colecciones de vacunas.
    \item \textbf{Vacuna:} mantiene nombre y fecha de aplicación o de próxima
    aplicación.
    \item \textbf{Fecha:} concentra validación, diferencia de días, cálculo de
    edad y comparación de mes.
    \item \textbf{Clínica:} administra el vector de pacientes, el archivo y el
    mutex que protege el estado compartido.
\end{itemize}

La separación evita que las reglas de fechas se repitan en las clases de
paciente y vacuna. También permite que las consultas temporales utilicen una
misma definición de edad y de intervalo, lo cual favorece la consistencia de
los resultados.

\subsection{Carga y validación}

Al construirse la clínica se obtiene la fecha del sistema y se carga el archivo
de pacientes. Cada línea se recorta, se divide en campos y se valida. Las
fechas se convierten a la representación `std::tm` solo después de comprobar
el número de día, el mes, el año y los años bisiestos. Las listas de vacunas se
procesan por pares; un par incompleto o inválido se omite sin impedir la carga
de los registros restantes.

Esta política privilegia la disponibilidad de la información válida: un error
local no invalida todo el archivo. Como contrapartida, el sistema no genera un
informe detallado de errores de entrada; esta limitación se registra en la
sección de discusión.

\subsection{Identificadores y orden}

Cada paciente recibe un identificador entero al construirse. Las búsquedas por
identificador utilizan búsqueda binaria, por lo que el vector se mantiene
ordenado por dicho identificador al insertar nuevos pacientes. Las búsquedas
por nombre y fecha recorren el vector porque esos atributos no poseen un índice
secundario.

El identificador es más apropiado que el índice como referencia concurrente:
una eliminación puede cambiar las posiciones del vector, pero el identificador
del paciente permanece estable durante la ejecución.

\subsection{Sincronización}

Las operaciones que acceden a la colección compartida adquieren el mutex de la
clínica. Esto incluye altas, búsquedas, modificaciones, eliminaciones y
consultas que recorren las listas. La salida de consola se protege además con
un mutex independiente para evitar que los mensajes de distintos hilos se
mezclen.

La regla aplicada es: toda lectura que deba observar un estado coherente y toda
modificación del vector se ejecutan bajo el mismo protocolo de exclusión. Las
funciones auxiliares que no toman el mutex solo se llaman desde una función que
ya lo posee; así se evita intentar adquirir dos veces un mutex no recursivo.

\subsection{Persistencia}

La clínica carga el archivo al iniciar y guarda el estado al destruirse. El
archivo se reconstruye a partir del vector actual, escribiendo una línea por
paciente y manteniendo la representación de sus vacunas. La reconstrucción
evita dejar estados parciales producto de múltiples ediciones y permite que la
siguiente ejecución utilice el archivo actualizado como punto de partida.

En términos de orden, los registros originales conservan su orden relativo
porque sus identificadores se asignan durante la carga; los nuevos registros
reciben identificadores posteriores y se insertan después de los originales.

% ---------------------------------------------------------------
\section{Diagrama de clases y grafo de dependencias}

\subsection{Diagrama de clases}

La Figura~\ref{fig:clases} muestra la estructura estática de la solución. Se
representan como clases las entidades principales y como componentes de utilidad
los archivos que concentran operaciones de fechas y comparación de texto. Este
diagrama no indica concurrencia: describe qué objetos se relacionan y qué
responsabilidad tiene cada uno.

\begin{figure}[H]
\centering
\begin{tikzpicture}[
    node distance=10mm and 12mm,
    uml/.style={draw, rectangle, align=left, text width=4.2cm, minimum height=1.8cm, font=\scriptsize, inner sep=5pt},
    utility/.style={draw, dashed, rectangle, align=center, text width=3.3cm, minimum height=1.0cm, font=\small},
    assoc/.style={-{Latex[length=2mm]}, thick},
    dep/.style={-{Latex[length=2mm]}, dashed, thick}
]
    \node[uml, fill=diagramgreen] (clinica) {
        \centering\textbf{Clinica}\\[-1mm]
        \rule{4.0cm}{0.3pt}\\[-1mm]
        \raggedright\texttt{- pacientes: vector<Paciente>}\\
        \texttt{- fechaActual: tm}\\
        \texttt{- mtx: mutex}\\[-1mm]
        \rule{4.0cm}{0.3pt}\\[-1mm]
        \texttt{+ agregarPaciente()}\\
        \texttt{+ buscar / modificar / consultar()}
    };
    \node[uml, fill=diagramorange, below left=of clinica] (paciente) {
        \centering\textbf{Paciente}\\[-1mm]
        \rule{4.0cm}{0.3pt}\\[-1mm]
        \raggedright\texttt{- id: int}\\
        \texttt{- nombre: string}\\
        \texttt{- fechaNacimiento: tm}\\
        \texttt{- vacunasRecibidas / próximas}\\[-1mm]
        \rule{4.0cm}{0.3pt}\\[-1mm]
        \texttt{+ agregar / modificar / eliminarVacuna()}
    };
    \node[uml, fill=diagramorange, below right=of clinica] (vacuna) {
        \centering\textbf{Vacuna}\\[-1mm]
        \rule{4.0cm}{0.3pt}\\[-1mm]
        \raggedright\texttt{- nombre: string}\\
        \texttt{- fechaAplicacion: tm}\\[-1mm]
        \rule{4.0cm}{0.3pt}\\[-1mm]
        \texttt{+ set / get / toString()}
    };
    \node[utility, fill=diagramblue, below=of paciente] (fecha) {\textbf{<<utility>> fecha.h}\\validar fecha, calcular edad\\y diferencia de días};
    \node[utility, fill=diagramblue, below=of vacuna] (utils) {\textbf{<<utility>> Utils}\\comparar texto\\y búsqueda binaria};
    \draw[assoc] (clinica) -- node[above left, font=\scriptsize] {contiene $0..*$} (paciente);
    \draw[assoc] (clinica) -- node[above right, font=\scriptsize] {administra} (vacuna);
    \draw[assoc] (paciente) -- node[above, font=\scriptsize] {contiene $0..*$} (vacuna);
    \draw[dep] (clinica) -- (fecha);
    \draw[dep] (clinica) -- (utils);
    \draw[dep] (paciente) -- (fecha);
    \draw[dep] (paciente) -- (utils);
    \draw[dep] (vacuna) -- (fecha);
\end{tikzpicture}
\caption{Diagrama de clases y componentes de utilidad.}
\label{fig:clases}
\end{figure}

`Clinica` es el punto de entrada de las operaciones y posee el vector compartido
de pacientes. `Paciente` compone sus vacunas, mientras que las operaciones de
fecha y texto se reutilizan desde componentes auxiliares. La política de
concurrencia se concentra en `Clinica`: las clases de dominio no necesitan
conocer qué hilo las está utilizando.

\subsection{Grafo de dependencias entre tareas y sección crítica}

El grafo de la Figura~\ref{fig:dependencias} representa las dependencias de
ejecución y el acceso de los hilos a la sección crítica. El recurso compartido
es el objeto `Clinica`, en particular su vector de pacientes. Las tareas verdes
de cada bloque pueden ejecutarse concurrentemente; sin embargo, las flechas
discontinuas muestran que todas deben atravesar la sección crítica protegida por
`mtx`. Por ello, las tareas son concurrentes en el nivel de planificación, pero
las modificaciones del vector se serializan una por una dentro del mutex.

\begin{figure}[H]
\centering
\begin{tikzpicture}[
    flow/.style={draw, rounded corners, align=center, minimum width=3.0cm, minimum height=0.72cm, font=\small},
    task/.style={draw, rounded corners, align=center, minimum width=2.35cm, minimum height=0.62cm, font=\scriptsize},
    resource/.style={draw, very thick, rounded corners, align=center, fill=red!15, minimum width=4.3cm, minimum height=0.9cm, font=\small},
    arr/.style={-{Latex[length=2mm]}, thick},
    access/.style={-{Latex[length=2mm]}, dashed, thick, gray!80!black}
]
    \node[flow, fill=diagramblue] (load) {Cargar y validar};
    \node[flow, fill=diagramblue, right=of load] (ids) {Preparar datos\\e identificadores};
    \node[task, fill=diagramgreen, below left=11mm and 11mm of ids] (a1) {T1 Alta Ana};
    \node[task, fill=diagramgreen, right=3mm of a1] (a2) {T2 Alta Beto};
    \node[task, fill=diagramgreen, right=3mm of a2] (a3) {T3 Alta Carla};
    \node[task, fill=diagramgreen, right=3mm of a3] (a4) {T4 Alta Diego};
    \node[task, fill=diagramgreen, right=3mm of a4] (a5) {T5 Alta Elena};
    \node[resource, below=10mm of a3] (cs1) {Sección crítica\\\texttt{Clinica.pacientes + mtx}};
    \node[flow, fill=diagramgreen, below=9mm of cs1] (join1) {COEND 1 / join};
    \node[task, fill=diagramgreen, below left=11mm and 12mm of join1] (b1) {T6 Vacunas Ana};
    \node[task, fill=diagramgreen, right=3mm of b1] (b2) {T7 Vacunas Beto};
    \node[task, fill=diagramgreen, right=3mm of b2] (b3) {T8 Próxima Carla};
    \node[task, fill=diagramgreen, below=4mm of b1] (b4) {T9 Próxima Diego};
    \node[task, fill=diagramgreen, right=3mm of b4] (b5) {T10 Modificar Ana/Beto};
    \node[task, fill=diagramgreen, right=3mm of b5] (b6) {T11 Eliminar Elena};
    \node[resource, below=10mm of b5] (cs2) {Misma sección crítica\\\texttt{Clinica.pacientes + mtx}};
    \node[flow, fill=diagramgreen, below=9mm of cs2] (join2) {COEND 2 / join};
    \node[flow, fill=diagramblue, below=of join2] (query) {Consultas y estadísticas};
    \node[flow, fill=diagramorange, right=of query] (save) {Guardar archivo};
    \draw[arr] (load) -- (ids);
    \draw[arr] (ids) -- (a3);
    \draw[access] (a1) -- (cs1);
    \draw[access] (a2) -- (cs1);
    \draw[access] (a3) -- (cs1);
    \draw[access] (a4) -- (cs1);
    \draw[access] (a5) -- (cs1);
    \draw[arr] (cs1) -- (join1);
    \draw[access] (b1) -- (cs2);
    \draw[access] (b2) -- (cs2);
    \draw[access] (b3) -- (cs2);
    \draw[access] (b4) -- (cs2);
    \draw[access] (b5) -- (cs2);
    \draw[access] (b6) -- (cs2);
    \draw[arr] (join1) -- (b2);
    \draw[arr] (cs2) -- (join2);
    \draw[arr] (join2) -- (query);
    \draw[arr] (query) -- (save);
\end{tikzpicture}
\caption{Grafo de dependencias y acceso concurrente a la sección crítica.}
\label{fig:dependencias}
\end{figure}

La primera dependencia es secuencial: se deben cargar los datos y preparar los
identificadores antes de iniciar las altas. T1--T5 forman un conjunto
concurrente. Después de `COEND 1`, T6--T11 también pueden crearse como hilos
independientes. El acceso a la sección crítica no elimina la concurrencia del
conjunto de tareas; únicamente impide que dos tareas modifiquen simultáneamente
la estructura `Clinica`. Finalmente, `COEND 2` es requisito para iniciar las
consultas y guardar un estado consistente.

% ---------------------------------------------------------------
\section{Representación COBEGIN--COEND}

La siguiente representación es pseudocódigo en español. No es el código fuente
del programa; describe la estructura de concurrencia que se implementa con
hilos y operaciones `join`.

\subsection{Alta concurrente}

\begin{lstlisting}[language={},caption={Pseudocódigo de altas concurrentes.},label={lst:altas}]
INICIO
    clínica <- cargar y validar pacientes.txt
    fecha_actual <- establecer fecha de trabajo

    COBEGIN
        tarea alta_Ana:
            si Ana no existe, agregar paciente Ana
        tarea alta_Beto:
            si Beto no existe, agregar paciente Beto
        tarea alta_Carla:
            si Carla no existe, agregar paciente Carla
        tarea alta_Diego:
            si Diego no existe, agregar paciente Diego
        tarea alta_Elena:
            si Elena no existe, agregar paciente Elena
    COEND

    conservar el identificador obtenido por cada tarea
FIN
\end{lstlisting}

Las cinco tareas consultan y actualizan el mismo objeto, por lo que la
independencia es lógica y no implica acceso simultáneo sin protección. Cada
alta adquiere el mutex antes de buscar e insertar. La inserción ordenada por
identificador mantiene la precondición de la búsqueda binaria.

\subsection{Vacunas, modificaciones y eliminación}

\begin{lstlisting}[language={},caption={Pseudocódigo de actualizaciones concurrentes.},label={lst:actualizaciones}]
COBEGIN
    tarea vacunas_Ana:
        registrar vacunas aplicadas para Ana

    tarea vacunas_Beto:
        registrar vacunas aplicadas para Beto

    tarea proxima_Carla:
        agregar vacuna próxima para Carla

    tarea proxima_Diego:
        agregar vacuna próxima para Diego

    tarea cambios_personales:
        modificar nombre de Ana
        modificar fecha de nacimiento de Beto

    tarea eliminar_Elena:
        eliminar a Elena si existe
COEND

realizar consultas solamente después del COEND
guardar el estado final al cerrar la clínica
\end{lstlisting}

Las tareas se refieren a pacientes distintos en sus operaciones principales.
Cuando dos tareas coinciden sobre el mismo paciente, como una aplicación de
vacuna y un cambio de datos personales, el mutex impide que se modifiquen
simultáneamente los campos del objeto. El resultado no depende de que una
operación se ejecute primero, siempre que ambas sean válidas y no exista una
dependencia funcional entre ellas.

\subsection{Correspondencia con la implementación}

En C++, `COBEGIN` se materializa creando un `std::thread` por tarea y `COEND`
se materializa invocando `join` sobre todos los hilos. El mutex de la clínica
protege el vector y el mutex de salida protege los mensajes de consola. De esta
forma, el pseudocódigo conserva la intención del modelo estructurado sin
introducir una sintaxis que el compilador de C++ no reconoce.

% ---------------------------------------------------------------
\section{Solución de cada requisito}

La Tabla~\ref{tab:requisitos} documenta cómo se resuelve cada parte de la
especificación y por qué se utiliza el mecanismo correspondiente.

\begin{table}[H]
\centering
\caption{Requisitos funcionales y solución implementada.}
\label{tab:requisitos}
\small
\begin{tabularx}{\textwidth}{p{0.7cm}p{4.0cm}YY}
\toprule
\textbf{No.} & \textbf{Requisito} & \textbf{Solución} & \textbf{Justificación} \\
\midrule
1 & Registrar paciente & Constructor de paciente y alta en la clínica. & Centraliza la validación y asigna un identificador estable. \\
2 & Registrar vacuna aplicada & Agregar a la lista de recibidas. & Distingue historial de vacunas futuras. \\
3 & Agregar vacuna próxima & Agregar a la lista de próximas. & Permite consultas temporales sin mezclar estados. \\
4 & Eliminar paciente & Búsqueda por identificador y borrado del vector. & El identificador evita depender de una posición que puede cambiar. \\
5 & Modificar nombre o nacimiento & Sobrecargas de modificación y validación de fecha. & Se reutiliza la misma identidad del paciente y se rechazan fechas imposibles. \\
6 & Mostrar información & Conversión del paciente a representación textual. & Proporciona una salida completa y legible. \\
7 & Modificar o eliminar vacuna & Búsqueda insensible a mayúsculas y operación sobre la lista. & Tolera variaciones de escritura en nombres de vacunas. \\
8 & Nombres en orden alfabético & Copia de nombres y ordenamiento por mayúsculas. & No altera el orden interno de persistencia. \\
9 & Mostrar información de cinco en cinco & Recorrido paginado con pausa cada cinco registros. & Evita una salida difícil de leer. \\
10 & Rango de edades & Cálculo de edad con fecha actual y filtrado. & La edad se calcula en años cumplidos, no solo por diferencia de años. \\
11 & Vacuna en próximos siete días & Diferencia de días y comparación inclusiva de 0 a 7. & Incluye tanto el día actual como el séptimo día futuro. \\
12 & Vacunas aplicadas en última semana & Recorrido de vacunas recibidas y conteo temporal. & Cuenta aplicaciones, no pacientes, que es lo solicitado. \\
13 & Vacuna del siguiente mes & Determinación del mes siguiente y filtrado por nombre. & Evita confundir una fecha del mismo mes con una del siguiente. \\
14 & Pacientes vacunados en mes actual & Filtrado, orden alfabético y salida de edad, vacuna y fecha. & Entrega todos los datos pedidos en una sola consulta. \\
\bottomrule
\end{tabularx}
\end{table}

% ---------------------------------------------------------------
\section{Pruebas y resultados}

\subsection{Compilación}

El reporte se organiza para compilarse desde su directorio mediante el comando
proporcionado para la práctica: `latexbuild main.tex`. El proceso utiliza
XeLaTeX, crea el directorio `build` y genera el PDF en `build/main.pdf`. Las
referencias bibliográficas se resuelven automáticamente por `latexmk` mediante
BibTeX y el estilo numérico IEEE.

El programa C++ se configura con CMake, requiere el estándar C++17 y enlaza la
biblioteca de hilos del sistema. El código fuente completo y el archivo de
configuración se incluyen en el Apéndice~\ref{app:codigo}, no en la explicación.

\subsection{Casos de prueba}

\begin{table}[H]
\centering
\caption{Casos de prueba considerados.}
\label{tab:pruebas}
\begin{tabularx}{\textwidth}{p{2.7cm}YY}
\toprule
\textbf{Caso} & \textbf{Procedimiento} & \textbf{Resultado esperado} \\
\midrule
Carga válida & Leer `pacientes.txt` con registros completos y campos vacíos permitidos. & Se crea un paciente por línea válida y se separan sus dos listas de vacunas. \\
Fecha inválida & Intentar crear una fecha con día fuera del mes o mes fuera de rango. & La validación rechaza la fecha mediante una excepción controlada. \\
Alta repetida & Ejecutar nuevamente el demostrador con pacientes ya existentes. & No se duplican los pacientes ni se acumulan las vacunas del demostrador. \\
Concurrencia & Ejecutar las altas y actualizaciones como tareas independientes. & Todos los hilos terminan y el vector permanece estructuralmente válido. \\
Búsquedas & Buscar por identificador, nombre y fecha de nacimiento. & Se obtiene el índice correspondiente o `-1` si no existe. \\
Consultas temporales & Consultar siete días, última semana, siguiente mes y mes actual. & Se filtran las vacunas según la fecha de trabajo establecida. \\
Persistencia & Cerrar la clínica y abrirla nuevamente. & Los cambios se reconstruyen desde el archivo actualizado. \\
\bottomrule
\end{tabularx}
\end{table}

La fecha de trabajo usada por el demostrador es el 13 de octubre de 2026. Una
fecha fija permite comparar ejecuciones y verificar resultados sin que cambien
por el paso del tiempo. En una aplicación operativa, dicha fecha puede
obtenerse del sistema o solicitarse al usuario, como ya contempla la clase de
la clínica.

\subsection{Resultados observados}

La compilación del programa en una copia limpia de los archivos fuente terminó
sin errores y produjo el ejecutable `vacunas`. La ejecución terminó con código
de retorno cero y no produjo mensajes en la salida de error. En la salida se
observaron cinco altas concurrentes, seis tareas de actualización y eliminación,
las búsquedas solicitadas y las siguientes estadísticas:

\begin{table}[H]
\centering
\caption{Resultados observados durante la ejecución de prueba.}
\label{tab:resultados}
\begin{tabularx}{0.78\textwidth}{p{7.0cm}Y}
\toprule
\textbf{Consulta} & \textbf{Resultado} \\
\midrule
Búsqueda por identificador de Ana & índice 0 \\
Búsqueda por nombre de Ana & índice 0 \\
Búsqueda por fecha 27/09/2020 & índice $-1$ \\
Vacunas aplicadas en la última semana & 2 \\
Vacunas de difteria en el siguiente mes & 1 \\
Estado de ejecución & código de retorno 0 \\
\bottomrule
\end{tabularx}
\end{table}

Estos resultados comprueban tanto el flujo concurrente como la continuación
secuencial posterior al `COEND`: las consultas se realizan cuando ya terminaron
las altas y modificaciones.

\subsection{Corrección de la sincronización}

La corrección se argumenta mediante tres propiedades:

\begin{enumerate}[leftmargin=1.2cm]
    \item \textbf{Exclusión mutua:} las operaciones protegidas por el mutex no
    se ejecutan simultáneamente sobre el vector.
    \item \textbf{Ausencia de consultas prematuras:} el `join` se ejecuta antes
    de las consultas, por lo que estas observan todas las actualizaciones
    terminadas.
    \item \textbf{Invariante de orden:} las inserciones conservan el orden por
    identificador, requisito necesario para la búsqueda binaria.
\end{enumerate}

La sincronización no hace que todas las operaciones sean paralelas: dos
operaciones que llegan al mismo mutex se serializan. El paralelismo se conserva
en la preparación y coordinación de tareas independientes, mientras que la
sección crítica se mantiene tan pequeña como permite la consistencia del
vector.

% ---------------------------------------------------------------
\section{Complejidad, limitaciones y discusión}

Sea $n$ el número de pacientes y $v$ el número total de vacunas almacenadas.
La carga del archivo es $O(n+v)$; una búsqueda binaria por identificador es
$O(\log n)$ en comparaciones, mientras que la búsqueda por nombre o fecha es
$O(n)$. Las consultas que recorren vacunas son $O(n+v)$ y los listados
ordenados requieren $O(n\log n)$ por el ordenamiento.

La implementación actual copia el vector en la función de búsqueda binaria,
por lo que el costo real de esa interfaz puede ser mayor que el costo ideal de
la búsqueda. Es una mejora futura pasar una referencia constante y evitar la
copia, sin modificar el algoritmo de búsqueda.

También existen las siguientes limitaciones:

\begin{itemize}[leftmargin=1.2cm]
    \item El formato de texto no ofrece transacciones ni bloqueo entre
    procesos distintos; la sincronización implementada protege los hilos del
    proceso actual.
    \item Los registros inválidos se ignoran durante la carga, pero no se
    reporta su número ni la causa exacta.
    \item Las utilidades de mayúsculas y minúsculas están orientadas a
    caracteres ASCII; la comparación de nombres con acentos puede requerir una
    normalización Unicode.
    \item Las funciones de salida mantienen el mutex mientras escriben; esto
    ofrece una vista consistente, pero puede reducir el paralelismo durante
    listados extensos.
    \item La escritura reconstruye el archivo completo. Para volúmenes grandes
    sería preferible una base de datos o un formato con actualizaciones
    transaccionales.
    \item El `CMakeLists.txt` utiliza una búsqueda recursiva de archivos fuente.
    Si se configura fuera del árbol fuente mientras existe un directorio de
    compilación generado dentro de él, podrían detectarse archivos `.cpp`
    auxiliares del compilador. Para una versión más robusta convendría declarar
    explícitamente las fuentes o excluir siempre los directorios de compilación.
\end{itemize}

Estas limitaciones no invalidan la solución de la práctica: delimitan el
alcance de una aplicación educativa de memoria compartida y muestran qué
decisiones deberían revisarse para un sistema de producción.

% ---------------------------------------------------------------
\section{Conclusiones}

El problema se resolvió separando el modelo de dominio, las reglas de fechas,
las operaciones de la clínica y la persistencia. Esta división facilita
justificar cada requisito y evita duplicar reglas de validación. El uso de
identificadores estables permite que las tareas concurrentes no dependan de
índices que pueden cambiar al insertar o eliminar pacientes.

El análisis de dependencias mostró que la carga y la asignación de
identificadores son precedencias obligatorias, mientras que las altas y varias
actualizaciones pueden agruparse en un bloque `COBEGIN`. El `COEND` se
representa con `join` y constituye la barrera que garantiza que las consultas
observen un estado terminado. El mutex evita condiciones de carrera sobre el
vector compartido, aunque también establece el límite del paralelismo efectivo.

Finalmente, la solución es adecuada para demostrar concurrencia estructurada,
validación, persistencia y consultas temporales en C++. Para una evolución
posterior se recomienda mejorar el diagnóstico de entradas inválidas, evitar
copias innecesarias en búsquedas, normalizar texto Unicode y utilizar un
almacenamiento transaccional cuando el volumen de datos lo requiera.

% ---------------------------------------------------------------
\clearpage
\addcontentsline{toc}{section}{Referencias}
\renewcommand{\refname}{Referencias}
\bibliographystyle{IEEEtran}
\bibliography{referencias}

% ---------------------------------------------------------------
\appendix
\section{Código fuente}
\label{app:codigo}

Este apéndice contiene el código fuente utilizado por la aplicación. Se omiten
los fragmentos de implementación en el cuerpo del reporte para mantener la
explicación centrada en el diseño, las dependencias y la justificación de las
decisiones.

\subsection{Configuración de compilación}
\lstinputlisting[language=C++,caption={Configuración CMake.}]{../programa/CMakeLists.txt}

\subsection{Programa principal}
\lstinputlisting[language=C++,caption={Programa principal.}]{../programa/main.cpp}

\subsection{Clínica}
\lstinputlisting[language=C++,caption={Interfaz de la clínica.}]{../programa/clinica.h}
\lstinputlisting[language=C++,caption={Implementación de la clínica.}]{../programa/clinica.cpp}

\subsection{Paciente y vacuna}
\lstinputlisting[language=C++,caption={Interfaz del paciente.}]{../programa/paciente.h}
\lstinputlisting[language=C++,caption={Implementación del paciente.}]{../programa/paciente.cpp}
\lstinputlisting[language=C++,caption={Interfaz de la vacuna.}]{../programa/vacuna.h}
\lstinputlisting[language=C++,caption={Implementación de la vacuna.}]{../programa/vacuna.cpp}

\subsection{Fechas y utilidades}
\lstinputlisting[language=C++,caption={Interfaz de fechas.}]{../programa/fecha.h}
\lstinputlisting[language=C++,caption={Implementación de fechas.}]{../programa/fecha.cpp}
\lstinputlisting[language=C++,caption={Interfaz de utilidades.}]{../programa/utils.h}
\lstinputlisting[language=C++,caption={Implementación de utilidades.}]{../programa/utils.cpp}

\section{Archivo de datos de prueba}
\label{app:datos}

El archivo utilizado para las pruebas se conserva como parte del material de
entrada de la aplicación.

\lstinputlisting[caption={Archivo de pacientes de prueba.}]{../programa/pacientes.txt}

\end{document}
